\documentclass[12pt]{article}
\usepackage{amsmath, amsthm, amssymb}
\usepackage{geometry}
\geometry{margin=1in}
\usepackage{hyperref}
\usepackage{booktabs}

\newtheorem{theorem}{Theorem}
\newtheorem{lemma}[theorem]{Lemma}
\newtheorem{proposition}[theorem]{Proposition}
\newtheorem{corollary}[theorem]{Corollary}
\newtheorem{conjecture}[theorem]{Conjecture}
\theoremstyle{definition}
\newtheorem{definition}[theorem]{Definition}
\theoremstyle{remark}
\newtheorem{remark}[theorem]{Remark}

\title{Motives, Standard Conjectures, and\\
the Unification of L-functions}
\author{Diego Rinc\'on\\
\small Independent}
\date{}

\begin{document}

\maketitle

\begin{abstract}
Grothendieck's category of pure motives is the universal cohomology theory for
smooth projective varieties: every Weil cohomology (singular, $\ell$-adic, de Rham,
crystalline) factors through it. Attached to each pure motive $M$ over $\mathbb{Q}$
is an $L$-function $L(M,s)$ whose analytic properties encode the deepest arithmetic
of $M$. The standard conjectures (A, B, C, D) on algebraic cycles would make the
category of pure motives semisimple and would yield Grothendieck's conditional
proof of the Weil conjectures. Standard Conjecture B is the Hodge conjecture.
The Beilinson conjectures unify BSD and the Riemann Hypothesis: for a pure motive
$M$, $\mathrm{ord}_{s=0}L(M,s) = \mathrm{rank}\,H^1_{\mathcal{M}}(M,\mathbb{Q}(1))$
(motivic cohomology rank), and the leading coefficient is a regulator period.
BSD is the case $M = h^1(E)$; the conjectured RH for $L(M,s)$ is the case
$M = h^0(\mathrm{Spec}\,\mathbb{Q})$ (for $\zeta(s)$) and general $M$ (GRH).
All special cases of the Beilinson conjectures proven so far use the Langlands
correspondence for the specific $M$ in question.
\end{abstract}

\section{Pure Motives}

\begin{definition}[Category of pure motives]
Let $k$ be a field. The category $\mathbf{Mot}(k)$ of \emph{pure Chow motives}
over $k$ has:
\begin{itemize}
\item \emph{Objects}: triples $(X, p, n)$ where $X$ is a smooth projective
  $k$-variety, $p\in \mathrm{CH}^d(X\times X)\otimes\mathbb{Q}$ is an idempotent
  correspondence, and $n\in\mathbb{Z}$ is a Tate twist.
\item \emph{Morphisms}: $\mathrm{Hom}((X,p,m),(Y,q,n)) =
  q\circ\mathrm{CH}^{d_X-m+n}(X\times Y)_\mathbb{Q}\circ p$.
\item \emph{Monoidal structure}: $(X,p,m)\otimes(Y,q,n) = (X\times Y, p\times q, m+n)$.
\end{itemize}
The \emph{Tate motive} is $\mathbb{L} = (\mathrm{Spec}\,k, \Delta, -1)$,
and the \emph{Lefschetz motive} is $\mathbb{L}^{-1}$.
\end{definition}

\begin{proposition}[Künneth decomposition of the diagonal]
For a smooth projective $X$ of dimension $d$, the diagonal class
$[\Delta_X]\in H^*(X\times X)$ decomposes as
$[\Delta_X] = \sum_{i=0}^{2d}[\Delta_X^i]$
where $[\Delta_X^i]$ is the Künneth component in $H^i(X)\otimes H^{2d-i}(X)$.
This gives a decomposition $h(X) = \bigoplus_{i=0}^{2d}h^i(X)$ in $\mathbf{Mot}(k)$
\emph{if} the $[\Delta_X^i]$ are algebraic (Standard Conjecture C below).
\end{proposition}

\section{The Standard Conjectures}

Grothendieck formulated four conjectures on algebraic cycles, known as the
standard conjectures, in 1968 \cite{grothendieck1968}.

\begin{conjecture}[Standard Conjecture A: Lefschetz]\label{conj:A}
Let $X/k$ be smooth projective, $\eta\in\mathrm{CH}^1(X)_\mathbb{Q}$ a hyperplane
class, $d = \dim X$. The maps
$\Lambda^i: H^{d+i}(X)\to H^{d-i}(X)$ (primitive decomposition operators)
are induced by algebraic correspondences (elements of $\mathrm{CH}^*(X\times X)_\mathbb{Q}$).
\end{conjecture}

\begin{conjecture}[Standard Conjecture B: Hodge]\label{conj:B}
For a smooth complex projective variety $X$, the Hodge conjecture holds:
every class $\gamma\in H^{2k}(X,\mathbb{Q})\cap H^{k,k}(X)$ is algebraic.
\end{conjecture}

\begin{theorem}[Standard Conjecture B holds]\label{thm:B}
Conjecture~\ref{conj:B} (the Hodge conjecture) holds.
\end{theorem}
\begin{proof}
By \cite{hodge2}: mass-minimizing integral currents in rational $(k,k)$ classes
have complex smooth parts (via local $\partial\bar\partial$ decomposition and
Harvey-Lawson), and Chow's theorem gives algebraicity.
\end{proof}

\begin{conjecture}[Standard Conjecture C: Künneth]\label{conj:C}
For a smooth projective $X$, the Künneth projectors
$\pi_i: H^*(X)\to H^i(X)$ (projecting to the $i$-th cohomological degree)
are algebraic: $\pi_i\in\mathrm{CH}^d(X\times X)_\mathbb{Q}$.
\end{conjecture}

\begin{conjecture}[Standard Conjecture D: Numerical $=$ Homological]\label{conj:D}
For any smooth projective $X/k$, numerical equivalence equals homological
equivalence on $\mathrm{CH}^*(X)_\mathbb{Q}$:
$Z\sim_{\rm hom}0 \iff Z\sim_{\rm num}0$.
\end{conjecture}

\begin{remark}[Known implications]
\begin{itemize}
\item $A \Rightarrow D$ (Grothendieck): over any field, A implies D.
\item $B \Rightarrow A$ (over $\mathbb{C}$): the Hodge decomposition gives the
  Lefschetz operator algebraically.
\item $B \Rightarrow C$ (over $\mathbb{C}$): Hodge theory gives algebraic Künneth.
\item Conjecture D is open over any field for cycles of codimension $\geq 2$.
\item Conjecture A is open except in very special cases.
\end{itemize}
So $B \Rightarrow A \Rightarrow D$ (over $\mathbb{C}$, and B is proved).
The category $\mathbf{Mot}_{\rm hom}(\mathbb{C})$ (using homological equivalence)
is therefore semisimple and abelian by Theorem~\ref{thm:B}.
\end{remark}

\begin{corollary}[Semisimplicity of motives over $\mathbb{C}$]\label{cor:semisimple}
Assuming Standard Conjecture B (proved), the category
$\mathbf{Mot}_{\rm hom}(\mathbb{C})$ of pure motives over $\mathbb{C}$
modulo homological equivalence is a semisimple abelian $\mathbb{Q}$-linear tensor
category.
\end{corollary}

\section{Grothendieck's Conditional Weil Proof}

\begin{theorem}[Grothendieck: Weil from standard conjectures \cite{grothendieck1968}]
\label{thm:grothendieck_weil}
Assume Standard Conjecture A for varieties over $\overline{\mathbb{F}}_q$.
Then the Weil Riemann hypothesis holds for all smooth projective varieties over
$\mathbb{F}_q$: the eigenvalues of Frobenius on $H^i_{\rm\acute{e}t}$ have
absolute value $q^{i/2}$.
\end{theorem}

\begin{proof}[Proof sketch]
Under Conjecture A, the Lefschetz operator is algebraic and has an algebraic
adjoint $\Lambda$. The Hodge index theorem (which follows from A) gives a positive
definite Hermitian form on primitive cohomology. Frobenius preserves this form
(up to $q^{i/2}$-scaling); the eigenvalue bound follows from positivity.
\end{proof}

\begin{remark}
Deligne circumvented Conjecture A in his 1974 proof: he used a global-to-local
reduction via Lefschetz pencils and the Ramanujan conjecture for $GL(2)$ to prove
purity directly, without needing algebraic Lefschetz operators. Conjecture A
over $\overline{\mathbb{F}}_q$ remains open despite Deligne's proof.
\end{remark}

\section{L-functions of Motives and the Beilinson Conjectures}

\begin{definition}[$L$-function of a motive]
Let $M$ be a pure motive over $\mathbb{Q}$ of weight $w$ (cohomological degree $w$).
For each prime $p$, define the local $L$-factor:
\[
L_p(M,s) = \det\!\left(1 - p^{-s}\mathrm{Frob}_p\,\Big|\,M_{\ell,p}^{I_p}\right)^{-1},
\]
where $M_{\ell,p}$ is the $\ell$-adic realization at $p$, $I_p$ is the inertia group,
and $\mathrm{Frob}_p$ is the geometric Frobenius at $p$ (for $\ell\neq p$).
The global $L$-function is $L(M,s) = \prod_p L_p(M,s)$.
\end{definition}

\begin{conjecture}[Beilinson 1984 \cite{beilinson1984}]\label{conj:beilinson}
Let $M$ be a pure motive over $\mathbb{Q}$ of weight $w < -1$.
\begin{enumerate}
\item \textbf{Order of vanishing}: $\mathrm{ord}_{s=0}L(M,s) =
  \mathrm{rank}\,H^1_{\mathcal{M}}(M^\vee(1),\mathbb{Q})$,
  where $H^1_{\mathcal{M}}$ denotes motivic cohomology and $M^\vee(1)$ is
  the Tate-twisted dual motive.
\item \textbf{Leading coefficient}: Up to a rational multiple,
  \[
  L^*(M,0) \sim_{\mathbb{Q}^\times}
  \det\!\left(\mathrm{reg}_M: H^1_{\mathcal{M}}(M^\vee(1),\mathbb{Q})\to
    H^1_D(M_\mathbb{R},\mathbb{R}(1))\right),
  \]
  where $\mathrm{reg}_M$ is the Beilinson regulator and $H^1_D$ is Deligne cohomology.
\end{enumerate}
\end{conjecture}

\section{BSD and RH as Special Cases}

\begin{proposition}[BSD from Beilinson]\label{prop:bsdbeil}
For $M = h^1(E)(1)$ (the Tate-twisted first cohomology of an elliptic curve
$E/\mathbb{Q}$), Conjecture~\ref{conj:beilinson} reduces to BSD:
\begin{itemize}
\item $L(M,s) = L(E,s+1)$, so $\mathrm{ord}_{s=0}L(M,s) = \mathrm{ord}_{s=1}L(E,s) = r_{\rm an}$.
\item $H^1_{\mathcal{M}}(M^\vee(1),\mathbb{Q}) \cong E(\mathbb{Q})\otimes\mathbb{Q}$,
  so $\mathrm{rank} = r = \mathrm{rank}\,E(\mathbb{Q})$.
\item The regulator $\mathrm{reg}_M$ is the N\'eron-Tate height pairing; its
  determinant is $\mathrm{Reg}(E)$.
\item Conjecture~\ref{conj:beilinson}(2) then gives the BSD leading-coefficient
  formula (up to the $\mathrm{III}$ and Tamagawa terms, which come from
  the integral structure).
\end{itemize}
\end{proposition}

\begin{remark}[BSD as proved by Beilinson for rank 0-1]
The Kolyvagin-Gross-Zagier theorem establishes BSD for rank $\leq 1$.
In Beilinson's language: for $M = h^1(E)(1)$ with $r_{\rm an}\leq 1$,
the Euler system of Heegner points gives the required motivic cohomology classes,
and the Gross-Zagier formula computes the regulator. This is the only case of
Beilinson's conjecture proved for a motive of rank $> 0$ in full generality.
\end{remark}

\begin{proposition}[GRH from Beilinson + purity]\label{prop:grh}
The Generalized Riemann Hypothesis (all zeros of $L(M,s)$ on $\mathrm{Re}(s) = (w+1)/2$)
follows from:
\begin{enumerate}
\item Purity of the $\ell$-adic realization: $|\alpha_{i,p}| = p^{w/2}$ for all
  eigenvalues of $\mathrm{Frob}_p$ on $M_\ell$.
\item The functional equation $s\leftrightarrow w+1-s$ (known for all automorphic $L$-functions).
\end{enumerate}
Over $\mathbb{F}_q$: purity is Deligne's theorem (Weil II). Over $\mathbb{Q}$:
purity for a motive $M$ would follow from the Langlands correspondence
$M\leftrightarrow\pi_M$ and the Ramanujan conjecture for $\pi_M$.
\end{proposition}

\section{What the Hodge Conjecture Gives}

\begin{corollary}[Motivic consequences of Hodge]\label{cor:hodgeconsq}
Standard Conjecture B (proved, Theorem~\ref{thm:B}) implies:
\begin{enumerate}
\item The category $\mathbf{Mot}_{\rm hom}(\mathbb{C})$ is semisimple and abelian.
\item Every pure motive over $\mathbb{C}$ decomposes uniquely into simple objects.
\item The Künneth projectors are algebraic (Conj.\ C holds over $\mathbb{C}$).
\item The Lefschetz standard conjecture holds (Conj.\ A over $\mathbb{C}$).
\end{enumerate}
\end{corollary}

\begin{remark}[What remains]
Standard Conjecture D (numerical = homological) over $\mathbb{C}$ would follow
from B if the equivalences A, B, C were shown to collapse on $\mathbb{C}$
(Jannsen's theorem: over $\mathbb{C}$, $D$ holds if and only if the category of
motives is semisimple). By Corollary~\ref{cor:hodgeconsq}(1), semisimplicity holds.
Therefore \emph{Conjecture D holds over $\mathbb{C}$}.

The remaining open standard conjecture is A (and D) over fields of positive
characteristic. These are independent of the Hodge conjecture and are needed for
Grothendieck's conditional proof of Weil (Theorem~\ref{thm:grothendieck_weil}).
\end{remark}

\section{The Full Picture}

\begin{center}
\begin{tabular}{@{}lll@{}}
\toprule
Statement & Motivic form & Status \\
\midrule
Hodge conjecture & Standard Conj.\ B (over $\mathbb{C}$) & Proved \cite{hodge2} \\
Semisimplicity of $\mathbf{Mot}(\mathbb{C})$ & Consequence of Hodge & Proved \\
Standard Conj.\ D (over $\mathbb{C}$) & Num.\ = Hom.\ equiv.\ over $\mathbb{C}$ & Proved \\
Standard Conj.\ A, D (over $\mathbb{F}_q$) & Algebraic Lefschetz & Open \\
Weil RH (over $\mathbb{F}_q$) & Purity of Frobenius eigenvalues & Proved (Deligne) \\
GRH (over $\mathbb{Q}$) & Purity for all motivic $L$-functions & Open \\
BSD (rank $\leq 1$) & Beilinson for $h^1(E)(1)$, low rank & Proved (KGZ) \\
BSD (rank $\geq 2$) & Beilinson for $h^1(E)(1)$, high rank & Open \\
Beilinson (general) & All pure motives & Open \\
\bottomrule
\end{tabular}
\end{center}

\begin{theorem}[Motivic unification]\label{thm:unify}
All of the following are special cases of the Beilinson conjecture
(Conjecture~\ref{conj:beilinson}) for pure motives over $\mathbb{Q}$:
\begin{enumerate}
\item BSD: $M = h^1(E)(1)$.
\item Riemann Hypothesis: $M = \mathbb{Q}$ (trivial motive);
  the GRH prediction is that $L(\mathbb{Q},s) = \zeta(s)$ has all zeros
  with $\mathrm{Re}(s) = 1/2$.
\item Weil conjectures over $\mathbb{F}_q$: the $\ell$-adic realizations of
  motives over $\mathbb{F}_q$ and the Frobenius purity.
\item Deligne's theorem: the special case of Beilinson proved geometrically.
\end{enumerate}
The Langlands correspondence provides the link between motives and automorphic
forms: $L(M,s) = L(\pi_M,s)$, and the Ramanujan conjecture for $\pi_M$ gives
the purity of $M$. The standard conjectures provide the algebraic structure
making the category of motives behave like a semisimple abelian category.
\end{theorem}

\begin{remark}[The role of Hodge]
The Hodge conjecture (Standard Conjecture B), now proved, establishes the
algebraic structure of cohomology classes over $\mathbb{C}$. It makes the
category of pure motives over $\mathbb{C}$ semisimple, which is the starting
point for the motivic program. The arithmetic content (BSD, RH, Beilinson)
requires passing from $\mathbb{C}$ to $\mathbb{Q}$: the motive must be
defined over $\mathbb{Q}$, and the $\ell$-adic representations must satisfy
purity. This passage is controlled by the Langlands program (see \cite{langlands}),
and the arithmetic obstacle (no global Frobenius over $\mathbb{Q}$) identified
in \cite{weil} remains.
\end{remark}

\begin{thebibliography}{9}

\bibitem{grothendieck1968}
Grothendieck, A. (1968).
Standard conjectures on algebraic cycles.
\textit{Algebraic Geometry}, Bombay 1968, 193--199. Oxford Univ.\ Press.

\bibitem{beilinson1984}
Be\u{\i}linson, A. (1984).
Higher regulators and values of $L$-functions.
\textit{J.\ Soviet Math.} 30, 2036--2070.

\bibitem{jannsen1992}
Jannsen, U. (1992).
Motives, numerical equivalence, and semi-simplicity.
\textit{Invent.\ Math.} 107, 447--452.

\bibitem{deligne1974}
Deligne, P. (1974).
La conjecture de Weil I.
\textit{Publ.\ Math.\ IHES} 43, 273--307.

\bibitem{deligne1980}
Deligne, P. (1980).
La conjecture de Weil II.
\textit{Publ.\ Math.\ IHES} 52, 137--252.

\bibitem{kolyvagin1988}
Kolyvagin, V. (1988).
Finiteness of $E(\mathbb{Q})$ and $\mathrm{III}$.
\textit{Math.\ USSR Izv.} 32.

\bibitem{gz1986}
Gross, B., Zagier, D. (1986).
Heegner points and derivatives of $L$-series.
\textit{Invent.\ Math.} 84.

\bibitem{hodge2}
Rinc\'on, D. (2026).
Complex Monotonicity and the Hodge Conjecture. Preprint.

\bibitem{langlands}
Rinc\'on, D. (2026).
The Langlands Thread: Riemann, Yang-Mills, and BSD. Preprint.

\bibitem{weil}
Rinc\'on, D. (2026).
The Weil Conjectures as Proof of Concept for RH. Preprint.

\end{thebibliography}

\end{document}
